%!TEX encoding = UTF8
%!TEX root = 0-notes.tex

\chapter{Problèmes de géométrie}
\label{chap:pb-geom}

Dans ce chapitre, et en mathématiques en général, les adjectifs \og orthogonal \fg \, et \og perpendiculaire \fg \, sont synonymes.
On parlera alors de projeté orthogonal d'un point sur une droite lorsque les droites engendrées sont perpendiculaires.

\section{Projeté orthogonal}

\exemulticols{1}{
	Un homme souhaite traverser une rivière en passant le moins de temps possible dans l'eau.
	Celui-ci part du point vert et dois arriver au point rouge.
	
	En supposant qu'il avance à vitesse constante dans l'eau, quel est le chemin idéal qu'il doit prendre ?
	\vfill\null
}{
	\null\vspace{-50pt}
	\begin{center}
	\includegraphics[page=23]{figures/fig-geom.pdf}
	\end{center}
}{exe:rivière}{
	Il existe plusieurs chemins (une infinité), leur point commun étant que leur section dans l'eau est orthogonale aux rives.
	En effet et d'après le cours, lorsque l'homme quitte la rive de départ, le point le plus proche sur l'autre rive est son projeté orthogonal sur celle-ci.
	La vitesse constante permet de lier la distance et le temps de façon proportionnelle.
}

\begin{multicols}{2}
	\begin{center}
	\includegraphics[page=8, scale=.9]{figures/fig-geom.pdf}
	\end{center}
	\columnbreak
	\dfn{projeté orthogonal}{
		Le \emphindex{projeté orthogonal} d'un point $M$ sur une droite $(d)$ est l'\emph{unique} point $M'$ tel que les droites $(MM')$ et $(d)$ soient perpendiculaires.
	}{}
\end{multicols}
	
\exe{}{
	Tracer une droite quelconque et un point n'appartenant pas à cette droite.
	À l'aide uniquement d'un compas, tracer le projeté orthogonal du point sur la droite.
}{exe:proj-compas}{
	TODO
}

\exemulticols{}{
	\begin{center}
	\includegraphics[page=17]{figures/fig-geom.pdf}
	\end{center}
}{
	\,\newline
	Dans la figure ci-contre, tracer un point $C$ tel que
		\begin{itemize}[rightmargin=-30pt]
			\item $A$ soit le projeté orthogonal de $C$ sur $(d)$ ; et 
			\item $B$ soit le projeté orthogonal de $C$ sur $(d')$.
		\end{itemize}
}{exe:proj-inv}{
	TODO
}
	
	% ça ne marche pas DU TOUT
	%\begin{vwcol}[widths={0.6,0.4},sep=.8cm, justify=flush,rule=10pt,indent=1em] 
	 %\end{vwcol}
\exemulticols{}{
	Dans la figure ci-contre, tracer
		\begin{enumerate}[rightmargin=-30pt]
			\item le point $B$, projeté orthogonal de $A$ sur $(d')$ ;
			\item le point $C$, projeté orthogonal de $B$ sur $(d)$ ; et
			\item le point $D$, projeté orthogonal de $C$ sur $(d')$.
		\end{enumerate}
	Que dire des droites $(AB)$ et $(CD)$ ? Justifier.
}{
	\begin{center}
	\includegraphics[page=18, scale=.8]{figures/fig-geom.pdf}
	\end{center}
}{exe:proj2}{
	\begin{center}
	\includegraphics[page=24, scale=1]{figures/fig-geom.pdf}
	\end{center}
	
	Par construction, les droites $(AB)$ et $(CD)$ sont toutes les deux perpendiculaires à la droite $(d')$.
	Elles sont donc parallèles.
}

\thm{distance minimale}{
	Le point $M'$, projeté orthogonal de $M$ sur $(d)$, est l'\emph{unique} point de $(d)$ le plus proche de $M$.
	
	Autrement dit, $M'$ est l'\emph{unique} minimiseur de la distance à $M$ en restant sur $(d)$.
}{thm:proj-min}



	\begin{center}
	\includegraphics[page=9]{figures/fig-geom.pdf}
	\end{center}

\pf%{Démonstration du théorème \ref{thm:proj-min}}
{}{
	Considérons un autre point $P \in (d)$ sur la droite $(d)$.
	Le triangle $MPM'$ est rectangle en $M'$, et le théorème de Pythagore implique donc que
		\[ MP^2 = MM'^2 + M'P^2. \]
	Or, comme un carré est toujours positif, on a les inégalités suivantes.
		\begin{align*}
			M'P^2 &\geq 0 \\
			M'P^2 + MM'^2 &\geq MM'^2 \\
			MP^2 &\geq MM'^2
		\end{align*}
	La distance de $M$ à $P$ est donc toujours plus grande ou égale à celle de $M$ à $M'$, ce qui démontre que $M'$ minimise sa distance à $M$ en restant sur $(d)$.
	
	Pour l'unicité, remarquons que l'égalité $MP^2 = MM'^2$ n'est vraie que lorsque $M'P^2 = 0$, c'est-à-dire lorsque $P=M'$. 
	$M'$ est donc bien l'unique point de $(d)$ le plus proche de $M$.
}


\exemulticols{, difficulty=3}{
	Considérons un cercle quelconque de centre $O$ et un point $P$ sur le cercle.
	Posons $(d)$ la droite, dite \emph{tangente au cercle}, touchant le cercle uniquement au point $P$.
	
	Démontrer que $(d)$ et $(OP)$ sont perpendiculaires.
	\vfill\null
}{
	\begin{center}
	\includegraphics[page=25, scale=1]{figures/fig-geom.pdf}
	\end{center}

}{exe:tangente-perp}{
	Fixons les objets $O, P, (OP), (d)$ de la figure et faisons varier le rayon du cercle.
	En partant d'un rayon nul (ou très petit) et en l'augmentant, on remarque que $P$ est le premier moment où le cercle et la droite $(d)$ se croisent. 
	En d'autres termes, le point $P$ est le point de $(d)$ le plus proche $O$ car un point plus proche appartiendrait à un cercle de rayon inférieur.
	D'après le cours, $P$ est le projeté orthogonal de $O$ sur $(d)$, et les droites $(d)$ et $(OP)$ sont perpendiculaires.
}

\exemulticols{}{
	Un homme est dans une salle pentagonale comme schématisé ci-contre.
	Celui-ci souhaite toucher les murs $[AB], [CD], [AE], [BC],$ et $[ED]$ dans cet ordre en effectuant le trajet le plus court à chaque déplacement.
	Tracer ses déplacements sur la figure si son point de départ est le point $M$.
}{
	\begin{center}
	\includegraphics[page=19]{figures/fig-geom.pdf}
	\end{center}	
}{exe:rayon-lumen}{
	TODO
}

\exe{, difficulty=1}{
	On considère une droite $(d)$ et les points $A$ et $B$ distincts n'appartenant pas à la droite $(d)$.
	On note $A'$ et $B'$ les projetés orthogonaux respectifs des points $A$ et $B$ sur la droite $(d)$.
	
	Calculer la longueur $AB$ sachant que $AA' =1, BB' = 5$, et $A'B' = 4$. 
	Distinguer deux cas selon que les points sont ou non du même côté de $(d)$.

}{exe:proj-d}{
	Lorsque $A$ et $B$ sont du même côté de $(d)$ on trouve $AB = 4\sqrt2$.
	
	Lorsque $A$ et $B$ sont de part et d'autre de $(d)$, on obtient $AB = 2\sqrt{13}$.
}

\ex{hauteur d'un triangle}{
	Pour calculer l'aire d'un triangle on crée un rectangle grâce à \emph{une} de ses hauteurs et, par symétrie, on déduit que
		\[ \text{Aire}(\text{triangle}) = \dfrac{\text{base} \cdot \text{hauteur}}2. \]
	
	\begin{center}
	\includegraphics[page=1, scale=1]{figures/fig-geom.pdf}
	\end{center}
	\warning\textbf{La hauteur est une fonction de la base. Si un autre côté est pris pour base, alors la hauteur sera différente.}
	La formule de l'aire ne dépend pas de la base choisie, mais il faut impérativement que la hauteur soit perpendiculaire à cette base, et qu'elle aille jusqu'au point culminant du triangle. 
	Ce point est le sommet du triangle qui n'appartient pas à la base.
	
	Autrement dit, la hauteur suit le projeté orthogonal du sommet n'appartenant pas à la base sur la droite engendrée par cette base.
	
}{}

\exemulticols{}{
	Soit $ABC$ le  triangle rectangle ci-contre, avec $BH$ une de ses hauteurs.
	\begin{enumerate}[label=\alph*.]
		\item Calculer $AB$.
		\item Calculer l'aire du triangle $ABC$.
		\item En déduire $BH$.
	\end{enumerate}
}{
	\begin{center}
	\includegraphics[page=16, scale=1.1]{figures/fig-geom.pdf}
	\end{center}
}{exe:aire-hauteur}{
	\begin{enumerate}[label=\alph*)]
		\item 
		Comme $AB^2 + BC^2 = AC^2$, on trouve
			\[ AB^2 = 104^2 - 96^2 = 1600 = 4\times100 = (2\times10)^2 = 20^2, \]
		d'où $AB = \pm20$ et donc $AB = 20$ par positivité.
		\item
		Le triangle étant rectangle, c'est la moitié d'un rectangle et son aire est donnée par
			\[ \frac12 AB \cdot BC = \frac12\times20\times96 = 960. \]
		\item En déduire $BH$.
		$BH$ est la hauteur du triangle $ABC$ de base $AC$.
		L'aire du triangle est donc donnée par
			\[ \frac12 AC \cdot BH = 960. \]
		Par conséquent, $BH = \frac{1920}{AC} = \frac{1920}{104} = \frac{240}{13} \approx 18,46$.
	\end{enumerate}
}


\exemulticols{, difficulty=1}{
	Considérons le triangle $ABC$, le point $D$ appartenant au segment $[AB]$.
	Supposons que $\text{Aire}(ABC) = \frac{3}{10}$ et que $AB=\frac34$ et $CD =\frac25$.
	
	Est-ce que les droites $(AB)$ et $(CD)$ sont perpendiculaires ?
	Raisonner par contradiction.
	\vfill\null
}{
	\begin{center}
	\includegraphics[page=26, scale=.6]{figures/fig-geom.pdf}
	\end{center}
} {exe:aire-to-perp}{
	Supposons que $(AB)$ et $(CD)$ soient perpendiculaires.
	Alors $D$ est le projeté orthogonal de $C$ sur $(AB)$ et $CD$ est la hauteur du triangle $ABC$ de base $AB$.
	L'aire du triangle $ABC$ est donc égale à
		\[ \frac12 CD\cdot AB = \frac12 \cdot \frac25\cdot\frac34 = \frac3{20}. \]
	Comme l'aire est égale à $\frac3{10}$ et non $\frac3{20}$, nous arrivons à une contradiction \Lightning.
	Les droites ne sont pas perpendiculaires.
}

\exe{, difficulty=2}{
	Considérons le triangle $ABC$ comme à l'exercice \ref{exe:aire-to-perp}, cette fois-ci d'aire $\text{Aire}(ABC) = 8$ et tel que $AB=4$ et $AC =5$.
	Quelle longueur $AD$ est nécessaire pour que les droites $(AB)$ et $(CD)$ soient perpendiculaires ?
} {exe:perp-to-aire}{
	Nous allons étudier ce que l'ajout de l'hypothèse d'orthogonalité des droites $(AB)$ et $(CD)$ implique sur $D$.
		\begin{align*}
			(AB) \perp (CD) && \implies && AD = ?
		\end{align*}
	L'orthogonalité des droites $(AB)$ et $(CD)$ implique que la hauteur du triangle $ABC$ de base $AB$ est égale à $CD$.
	Or, connaissant l'aire et la base, une seule hauteur est possible.
	Dès lors, si $(AB)$ et $(CD)$ sont perpendiculaires, alors
		\[ \frac12AB \cdot CD = 8 \iff CD = \frac{16}{4} = 4. \]
		
	En outre, si le triangle $ACD$ est rectangle en $D$, alors $AC^2 = AD^2 + CD^2$.
	Ceci implique que
		\[ AD^2 = AC^2 - CD^2 = 5^2 - 4^2 = 9, \]
	et donc que $AD = 3$, car la longueur $AD$ est positive.
	
	Le seul choix possible est donc $AD = 3$ pour que $(AB)$ et $(CD)$ soient perpendiculaires.
}

\section{Trigonométrie}

Les angles sont généralement dénotés par des lettres grecques minuscules.
	\begin{align*}
		\alpha &: \text{\og alpha \fg} & \beta &: \text{\og beta \fg} & \gamma &: \text{\og gamma \fg} \\
		\theta &: \text{\og theta \fg} &  \delta &: \text{\og delta \fg} & \omega &: \text{\og omega \fg}
	\end{align*}



\subsection{Théorème de Thalès}

On rappelle les formules de trigonométrie uniquement applicables dans un triangle rectangle.
\warning\textbf{Les côtés opposé et adjacent dépendent de l'angle choisi.}


\dfn{}{
	\begin{align*}
		\text{sinus} = \dfrac{\text{côté opposé}}{\text{hypoténuse}}
		&& 
		\text{cosinus} = \dfrac{\text{côté adjacent}}{\text{hypoténuse}}
		&&
		\text{tangente} = \dfrac{\text{côté opposé}}{\text{côté adjacent}}
	\end{align*}

	\begin{center}
	\includegraphics[page=10, scale=1]{figures/fig-geom.pdf}
	\hspace{60pt}
	\includegraphics[page=11, scale=1]{figures/fig-geom.pdf}
	\end{center}
Moyens mnémotechniques : 
	\begin{center}
	SOH CAH TOA ou CAH SOH TOA.
	\end{center}

}{dfn:sin-cos-tan}

% à supprimer ? on a le quart de cercle trigonométrique de toute façon
Ces formules sont en fait une conséquence du théorème de Thalès : en connaissant deux angles et un côté d'un triangle, on peut l'agrandir ou le rapetisser pour maintenir les proportions et les angles et se ramener à un triangle rectangle étalon d'hypoténuse de longueur $1$.

	\begin{center}
	\includegraphics[page=2, scale=1]{figures/fig-geom.pdf}
	\end{center}

En effet, les deux triangles ci-dessus sont semblables (on obtient l'un en (dé)zoomant l'autre), et en \emph{définissant} les longueur des côtés du deuxième par $\cos(\alpha)$ et $\sin(\alpha)$, le théorème de Thalès nous donne, par exemple,
	\[ \dfrac{\cos(\alpha)}{1} = \dfrac{\text{adjacent}}{\text{hypoténuse}}. \]

En conclusion, le cosinus, le sinus, et la tangente sont donc en fait une banque de données de longueurs de côtés d'un triangle particulier qu'on a mesuré à l'avance.
Certaines formules mathématiques permettent de calculer exactement ces valeurs, mais en général elle sont approximatives (mais aussi précises qu'on le souhaite !) et il n'existe pas de formule compacte pour les exprimer.

Remarquons immédiatement qu'on a la relation
	\begin{align*}
		\dfrac{\text{sinus}}{\text{cosinus}}
		= \dfrac{\text{opposé}}{\text{hypoténuse}} \cdot \dfrac{\text{hypoténuse}}{\text{adjacent}}
		= \dfrac{\text{opposé}}{\text{adjacent}}
		 = \text{tangente}.
	\end{align*}


\thm{}{
	Soit $\theta \in [0; 90\degree]$ un angle aigu ou droit.
	Alors 
		\[ \tan(\theta) = \dfrac{\sin(\theta)}{\cos(\theta)}. \]
}{}

\exemulticols{, difficulty=1}{
	Soit un triangle $ABC$ quelconque représenté ci-contre. On pose $a=a_1 + a_2 = BC$.
	
	Montrer que
		$\text{Aire}(ABC) = \dfrac12 ab \sin(\gamma). $
}{
	\begin{center}
	\includegraphics[page=20, scale=.7]{figures/fig-geom.pdf}
	\end{center}
}{exe:formule-aire}{
	L'aire $ABC$ est égale à $\frac12 h \cdot a$.
	Or, $\sin(\gamma) = \frac{h}{b} \iff h = b\sin(\gamma)$.
	En substituant l'expression de $h$ dans celle de l'aire, on conclut.
}

\exe{, difficulty=2}{
	On souhaite démontrer le \emphindex{théorème d'Al-Kashi}\footnotemark, ou \emphindex{loi des cosinus}, généralisant le théorème de Pythagore aux triangles non nécessairement rectangles.
	Reconsidérons donc le triangle $ABC$ de l'exercice \ref{exe:formule-aire}.
	
	\begin{enumerate}
		\item Montrer d'abord que $a_2 = b \cos(\gamma)$.
			%\begin{align*} h = b \sin(\gamma), && \text{ et } && a_2 = b \cos(\gamma). \end{align*}
		\item
		Donner les identités de Pythagore dans les deux triangles rectangles $HAC$ et $BAH$, où $H$ est le projeté orthogonal de $A$ sur $(BC)$.
		\item
		Utiliser $a_1 = a - a_2$ et les expressions obtenues pour conclure que
			\[ c^2 = a^2 + b^2 - 2ab \cos(\gamma). \]
	\end{enumerate}
}{exe:al-kashi}{
	Posons $H$ le projeté orthogonal de $A$ sur $(BC)$.
	\begin{enumerate}
		\item 
			Le résultat est obtenu par trigonométrie dans le triangle rectangle $HAC$.
		\item
		On obtient
			\begin{align*}
				c^2 &= h^2 + a_1^2, \\
				b^2 &= h^2 + a_2^2. 
			\end{align*}
		\item
		En partant de la gauche on a 
			\begin{align*}
				c^2 &= h^2 + a_1^2, \\
					&= h^2 + (a-a_2)^2, \\
					&= h^2 + a_2^2 + a^2 - 2aa_2, \\
					&= b^2 + a^2 - 2ab\cos(\gamma),
			\end{align*}
		comme requis.
	\end{enumerate}
}

\footnotetext{Ghiyath ad-Din Jamshid Mas`ud al-Kashi (1380? - 1429), mathématicien et astronome perse.}

\exe{, difficulty=1}{
	Soit un triangle aux angles $\alpha, \beta, \gamma$ opposés au côtés de longueur $a, b, c$ respectivement.
	En considérant d'abord la hauteur $h$ partageant l'angle $\gamma$, montrer que $h = a\sin\beta = b\sin\alpha$.
	
	En déduire la \emphindex{loi des sinus} :
		\[ \dfrac{a}{\sin\alpha} = \dfrac{b}{\sin\beta} = \dfrac{c}{\sin\gamma}. \]
}{exe:loi-sinus}{
	La trigonométrie simple donne les égalités $h = a\sin\beta = b\sin\alpha$, desquelles on déduit que 
		\[ \dfrac{a}{\sin\alpha} = \dfrac{b}{\sin\beta}. \]
	En considérant une autre hauteur partageant un autre angle, on obtient la dernière égalité de la loi des sinus.
}

\subsection{Quart de cercle trigonométrique}
\label{sec:quart-trigo}

\begin{multicols}{2}
Considérons, dans le quadrant supérieur-droit d'un repère, un quart de cercle de rayon $1$ centré en l'origine $O$.

On choisit un angle $\theta \in [0, 90\degree]$, et on nomme le point $P$ du quart de cercle tel que l'angle entre $(OP)$ et l'axe des abscisse soit $\theta$.

Le triangle rectangle ainsi créé est d'hypoténuse $1$ et vérifie donc les propriétés du triangle étalon considéré dans la section précédente.
On a donc, par définition,
	\begin{align*}
		x = \cos(\theta), && \text{ et } && y = \sin(\theta).
	\end{align*}

	\begin{center}
	\includegraphics[page=3, scale=1]{figures/fig-geom.pdf}
	\end{center}
	
\end{multicols}

Si on a pas peur des raisonnements circulaires, on peut revérifier les égalités avec les formules apprises.
	\begin{align*}
		\cos(\theta) = \dfrac{\text{adjacent}}{\text{hypoténuse}} = \dfrac{x}{1} = x, 
		&& \et &&
		\sin(\theta) = \dfrac{\text{côté opposé}}{\text{hypoténuse}} = \dfrac{y}1 = y.
	\end{align*}
En conclusion, les coordonnées du point $P$ sont exactement
	\[ P = (x ; y) = \bigl( \cos(\theta) ; \sin(\theta) \bigr). \]
Dans ce contexte, on déduit naturellement le théorème suivant.

\notations{
	Comme il est assez fastidieux d'écrire $ \bigl(\cos\left(\theta\right)\bigr)^2$,
	on optera pour la notation plus allégée suivante :
		\[\cos^2 (\theta) =  \bigl(\cos\left(\theta\right)\bigr)^2. \]
}{}

\thm{}{
	Soit $\theta \in [0\degree; 90\degree]$ un angle aigu ou droit.
	Alors, 
		\begin{align*}
			0 \leq \cos(\theta) \leq 1, &&
			0 \leq \sin(\theta) \leq 1,
		\end{align*}
	avec $\sin(0\degree) = 0, \sin(90\degree) = 1$ et $\cos(0\degree) = 1, \cos(90\degree) = 0$.
	De plus,
		\begin{gather*}
			\cos^2(\theta) + \sin^2(\theta) = 1.
		\end{gather*}
}{thm:c2-s2}

\exemulticols{}{
	On considère un point $P$ sur un quart de cercle de rayon $2$ et de centre $O$, l'origine d'un repère.
	On pose $\theta \in [0; 90\degree]$ l'angle entre la droite $(OP)$ et l'axe des abscisses.
	Chaque choix de $\theta$ correspond donc à un point et aussi à un triangle comme représentés ci-contre.
	
	\begin{enumerate}
		\item
		Montrer que l'aire $f(\theta)$ du triangle est donnée par 
			\[ f(\theta) = 2 \cos(\theta) \sin(\theta). \]
		\item
		Grapher $\C_f$ à l'aide de la calculatrice et en déduire l'angle maximisant $f$ ainsi que les coordonnées du point $P$ correspondant.
	\end{enumerate}
}{
	\begin{center}
	\includegraphics[page=21]{figures/fig-geom.pdf}
	\end{center}
}{exe:opt-area}{
	\begin{enumerate}
		\item
		Par trigonométrie, $x = 2\cos(\theta)$ et $y = 2\sin(\theta)$.
		Le triangle étant rectangle, son aire est la moitié de l'aire du rectangle de côtés $x$ et $y$, égale à
			\[ \dfrac12 xy = \frac12 2\cos(\theta)2\sin(\theta) = f(\theta), \]
		comme requis.
		
		\item
		\begin{center}
		\includegraphics[page=27]{figures/fig-geom.pdf}
		\end{center}
		On lit $x^\star \approx 45^\circ$ (c'est en fait une égalité) et $f(x^\star) \approx 1$.
		Le point $P$ correspondant a pour coordonnées 
			\[ P\bigl(2\cos(45^\circ) ; 2\sin(45^\circ)\bigr) = \bigl(\sqrt2 ; \sqrt2\bigr). \]
	\end{enumerate}
}

\hrule

\exemulticols{}{
	On considère exactement la même situation qu'à l'exercice précédent mais cette fois ci en prenant $x$ l'abscisse du point $P$ du quart de cercle comme paramètre.
	On a donc $x \in [0; 2]$.
	
	\begin{enumerate}
		\item Montrer que l'aire du rectangle est donnée par
			\[ g(x) = x \sqrt{4-x^2}. \]
		\item
			Estimer la valeur de $x$ qui maximise $g$.
		\item
			Déduire de l'exercice précédent que $f(x)$ est maximal en $x=\sqrt{2}$.
	\end{enumerate}
}{
	\begin{center}
	\includegraphics[page=22]{figures/fig-geom.pdf}
	\end{center}
}{exe:cossin2}{
	TODO
}


\section{Calculs exacts}

\subsection{Calculs exacts de $\cos(\theta), \sin(\theta), \tan(\theta)$ pour $\theta=30^\circ, 60^\circ$}

Considérons le triangle équilatéral suivant.
Par symétrie, tous les angles sont égaux à 60°.
En prenant une hauteur quelconque, on obtient un triangle rectangle 30-60-90 comme suit.

	\begin{center}
	\includegraphics[page=12]{figures/fig-geom.pdf}
	\hspace{4cm}
	\includegraphics[page=13]{figures/fig-geom.pdf}
	\end{center}

Le théorème de Pythagore donne la longueur $h$ de la hauteur : 
	\[ 1^2 = \left(\dfrac12\right)^2 + h^2 \iff h^2 = \dfrac34 \iff h = \dfrac{\sqrt3}2, \]
On peut donc connaître les valeurs de $\cos(\theta), \sin(\theta), \tan(\theta)$ pour $\theta\in\bigset{ 30^\circ ; 60^\circ}$.

\mprop{}{
	% moche de faire comme ça mais ça donne le meilleur rendu imo
		\[ \sin(30^\circ) = \dfrac12, \qquad \cos(30^\circ) = \dfrac{\sqrt{3}}2, \qquad \tan(30^\circ) = \dfrac{\sqrt{3}}3. \]
	
		\[ \cos(60^\circ) = \dfrac12, \qquad \sin(60^\circ) = \dfrac{\sqrt{3}}2, \qquad \tan(60^\circ) =\sqrt{3}.\]
}{prop:données}

\exe{, difficulty=1}{
	Calculer $\cos^2(\theta) + \sin^2(\theta)$ pour $\theta = 30^\circ$ et $\theta =60^\circ$ à l'aide des valeurs exactes données dans le cours.
}{exe:c2s2-vérif}{
	Pour $\theta = 30^\circ$, le cours donne $\cos(30^\circ) = \frac{\sqrt3}2$ et $\sin(30^\circ) = \frac12$.
	En sommant leur carré, on trouve
		\begin{align*}
			\cos^2(30^\circ) + \sin^2(30^\circ) &= \left(\frac{\sqrt3}2\right)^2 + \left(\frac12\right)^2 \\
			&= \frac34 + \frac14 \\
			&= 1
		\end{align*}
	Un calcul analogue donne $\cos^2(60^\circ) + \sin^2(60^\circ) = 1$ également.
}


\exe{, difficulty=1}{
	Vérifier qu'on a bien $\tan(\theta) = \dfrac{\sin(\theta)}{\cos(\theta)}$ pour $\theta = 30^\circ, 60^\circ$ à l'aide des valeurs exactes données dans le cours.
}{exe:tan-vérif}{
	Pour $\theta = 30^\circ$, le cours donne $\cos(30^\circ) = \frac{\sqrt3}2, \sin(30^\circ) = \frac12,$ et $\tan(30^\circ) = \frac1{\sqrt3}$.
	
	Comme diviser équivaut à multiplier par l'inverse, il suit que
		\begin{align*}
			\dfrac{\sin(30^\circ)}{\cos(30^\circ)} &= \dfrac{\frac12}{\frac{\sqrt3}2} \\
			&= \frac12 \cdot \frac2{\sqrt3} \\
			&= \frac1{\sqrt3} \\
			&= \tan(30^\circ).
		\end{align*}
	Un calcul analogue permet de vérifier l'identité pour $\theta = 60^\circ$.
}

\subsection{Calculs exacts de longueurs : manipulation des racines carrées}

Le calcul exact de longueurs implique de pouvoir manipuler les racines carrées dans les équations.
La racine et ses propriétés ont été étudiées dans la section \ref{sec:racine-carrée} du chapitre \ref{chap:fonction-carré}.

\ex{}{
	Pour résoudre pour $x\in\R$, $x\neq0$, l'équation 
		\[ \dfrac8x  =\dfrac{\sqrt{2}}2, \]
	on souhaite ramener $x$ au numérateur afin d'obtenir quelque chose de la forme $x = \dots$.
	Pour ça, on multiplie par $x$ les deux membres de l'égalité pour trouver
		\[ 8 = \dfrac{\sqrt{2}}2 x. \]
	On multiplie ensuite par $2$ et par $\frac1{\sqrt{2}}$ pour obtenir ce qu'on voulait :
		\begin{align*}
			8 \cdot 2 \cdot \dfrac1{\sqrt{2}} = x, && \iff && x = \dfrac{16}{\sqrt{2}}.
		\end{align*}
	On peut manipuler davantage l'expression en multipliant le numérateur et le dénominateur par $\sqrt{2}$ pour trouver
		\[ x = \dfrac{16 \sqrt{2}}{\sqrt{2}^2} = 8 \sqrt{2}. \]
	Les deux nombres $\frac{16}{\sqrt{2}}$ et $8 \sqrt{2}$ sont égaux, au même titre que $\frac36 = \frac12$.
	La deuxième expression est cependant plus agréable à lire et à écrire.
}{ex:eq-sqrt}

\qs{}{
	Peut-on simplifier encore davantage l'expression $8 \sqrt{2}$ ?
	Autrement dit, peut-on écrire $\sqrt2$ comme fraction de deux entier à multiplier par $8$ ?
	$\sqrt{2}$ est-il un nombre rationnel ou irrationnel ?
}

\thm{}{
	Soit $d \in\N$ un entier naturel qui n'est pas un carré parfait ($d\neq k^2, k\in\N$).
	Alors $\sqrt{d}$ n'est pas rationnel.
}{thm:sqrt-irrational}

\cor{}{
	La racine carrée de $2$ est irrationnelle : $\sqrt2 \notin \Q.$
	Son écriture décimale est infinie et apériodique.
}{cor:sqrt2-irrational}

\pf{démonstration du corrolaire \ref{cor:sqrt2-irrational}}{
	Par l'absurde, supposons que $\sqrt2 \in \Q$ soit rationnel.
	
	Par définition, $\sqrt2 = \frac{a}{b}$ avec $a, b \in \N$ deux entiers.
	Ainsi, par mise au carré,
		\begin{align*}
			2 = \left(\dfrac{a}b\right)^2 = \dfrac{a^2}{b^2} && \iff && a^2 = 2b^2.
		\end{align*}
	Étudions cette égalité d'entiers à l'aide du corollaire \ref{cor:carré-paire}.
	À gauche, $a^2$ est un carré parfait, et donc la puissance de $2$ dans sa décomposition en facteurs premiers est nécessairement paire.
	À droite, la même chose tient pour $b^2$. Il suit que, dans la décomposition en facteurs premiers de $2b^2$, la puissance de $2$ est impaire.
	
	L'entier $a^2$ admet donc deux décompositions en facteurs premiers différentes, ce qui contredit le théorème fondamental de l'arithmétique (théorème \ref{thm:decomposition-primaire-strong}). {\Large\Lightning}
	
	Concernant l'écriture décimale d'un nombre irrationnel, voir le théorème \ref{thm:décimal-périodique}.
}

\nt{
	La démonstration ci-dessus s'appuie sur l'unicité de la décomposition primaire du théorème \ref{thm:decomposition-primaire-strong}.
	Ce résultat étant assez difficile à prouver, la démonstration de l'irrationnalité de $\sqrt2$ est parfois donnée de la façon suivante, sans utiliser le théorème fondamental de l'arithmétique
}

\pf{démonstration du corrolaire \ref{cor:sqrt2-irrational} sans théorème fondamental de l'arithmétique}{
	Reprenons la démonstration ci-dessus en supposant, de surcroît, que la fraction $\frac{a}b$ est irréductible.
	
	À partir de l'identité $a^2 = 2b^2$, on déduit que $a^2$ est pair.
	Il suit que
		\[ a^2 - 1 = (a+1)(a-1) \]
	est impair.
	Par conséquent, $a+1$ et $a-1$ sont tous les deux impairs, et donc $a$ est pair.
	Nous avons donc montré
		\begin{align*}
			a^2 \text{ pair} && \implies && a \text{ pair}.
		\end{align*}
	Il suit que $a = 2k$ pour un entier $k\in\Z$, et donc que $(2k)^2 = 2b^2 \iff b^2 = 2k^2$.
	
	Les rôles sont à présents inversés : c'est $b^2$ qui est pair. 
	Le même raisonnement mène à ce que $b$ soit pair, une contradiction car la fraction $\frac{a}b$ a été supposée irréductible. {\Large\Lightning}
}


\exe{, difficulty=2}{
	Démontrer le théorème \ref{thm:sqrt-irrational}.
}{exe:sqrt-irrational}{
	Si $\sqrt{d} = \dfrac{a}{b}$, alors $a^2 = d\cdot b^2$.
	Comme $d$ n'est pas un carré parfait, un des premiers de sa décomposition apparaît à puissance impaire.
	Ce premier apparaît à puissance paire dans $b^2$ et donc impaire dans $a^2$, ce qui est une contradiction. \Large\Lightning
}

\exe{, difficulty=2}{
	Soit $x\in\R$ un nombre vérifiant $x^3 = 2$.
	Montrer que $x$ n'est pas rationnel : $x\not\in\Q$.
}{exe:cbrt2-irr}{
	Supposons que $x = \dfrac{a}b$. En mettant au cube, on obtient $a^3 = 2b^3$, égalité d'entiers naturels. 
	Par unicité de la décomposition de facteurs premiers, les deux nombres admettent la même décomposition.
	À gauche, le nombre premier 2 apparaît à une puissance qui est un multiple de 3.
	À droite, le premier 2 apparaît à une puissance qui est un multiple de 3 plus 1. \Large\Lightning
}

\exe{, difficulty=2}{
	En musique, la gamme classique contient douze notes avec un rapport de fréquence de 2 sur l'octave.
	La gamme est dite à tempérament égal si le rapport de fréquence entre deux notes consécutives est toujours le même.
	Notons ce rapport $r\in\R$.
	Justifier que $r^{12} = 2$ puis démontrer que $r$ est un nombre irrationnel : $r \not\in\Q$.
}{exe:temp-égal}{
	Le rapport de fréquence entre deux notes consécutives étant toujours $r$, on multiplie $r$ par lui-même 12 fois pour obtenir le rapport de fréquence entre deux notes de distance $12$, c'est-à-dire deux notes séparée par un octave.
	Ainsi $r^{12} = 2$ car l'octave correspond à un rapport de fréquence de 2.
	
	Supposons désormais que $r$ soit rationnel. Alors $r^{6}$ est aussi rationnel.
	Cependant, $(r^6)^2 = r^{12} = 2$, donc $r^{6} = \sqrt2$, qui n'est pas rationnel. \Large\Lightning
}

\mprop{manipulation des racines carrées}{
	Soit $d\in\R_+$ un nombre réel positif qui n'est pas un carré parfait.
	Alors
		\begin{align*}
			\dfrac1{\sqrt{d}} = \dfrac{\sqrt{d}}d, && \et && \dfrac1{1+\sqrt{d}} = \dfrac{1-\sqrt{d}}{1-d}.
		\end{align*}
}{prop:conjugué}
\pf{}{
	On utilise l'identité remarquable $(a+b)(a-b) = a^2 - b^2$ pour supprimer la racine du dénominateur.
	\begin{align*}
		\dfrac1{1+\sqrt{d}} &= \dfrac{1-\sqrt{d}}{(1+\sqrt{d})(1-\sqrt{d})}
							= \dfrac{1-\sqrt{d}}{1-d}.
	\end{align*}
}

\nomen{
	On appelle $a - b \sqrt{d}$ l'\emphindex{expression conjuguée} de $a + b \sqrt{d}$ (et vice versa).
}



\exe{, difficulty=1}{
	Exprimer les fractions suivantes sous la forme $p + q\sqrt{d}$ avec $p, q\in\Q$ rationnels et $d \in\N$ entier naturel.
	\begin{multicols}{3}
	\begin{enumerate}
		\item $\dfrac1{\sqrt5}$
		\item $\dfrac5{2\sqrt7}$
		\item $\dfrac1{1+\sqrt2}$
		\item $\dfrac{2+\sqrt2}{\sqrt2+3}$
		\item $\dfrac{2\sqrt3}{6-\sqrt3}$
		\item $\dfrac{\sqrt7+6}{1+5\sqrt7}$
	\end{enumerate}
	\end{multicols}
}{exe:conjugue}{
	\begin{multicols}{2}
	\begin{enumerate}
		\item 
			\begin{align*}
				\dfrac1{\sqrt5} = \dfrac{\sqrt5}5 = \frac15 \sqrt5
			\end{align*}
		\item
			\begin{align*}
				\dfrac5{2\sqrt7} = \dfrac5{14}\sqrt7
			\end{align*}
		\item 
			\begin{align*}
				\dfrac1{1+\sqrt2} &= \dfrac{\sqrt2-1}{(1+\sqrt2)(\sqrt2-1)} \\
				&= \dfrac{\sqrt2-1}{2-1} \\
				&= \sqrt2 -1
			\end{align*}
		\item 
			\begin{align*}
				\dfrac{2+\sqrt2}{\sqrt2+3} &= \dfrac{(2+\sqrt2)(3-\sqrt2)}{(\sqrt2+3)(3-\sqrt2)} \\
				&= \dfrac{6+2\sqrt2+3\sqrt2 - 2}{9-2} \\
				&= \frac{4+5\sqrt2}{7} \\
				&= \frac47 + \frac57\sqrt2
			\end{align*}
		\item 
			\begin{align*}
				\dfrac{2\sqrt3}{6-\sqrt3} &= \dfrac{2\sqrt3 (6+\sqrt3)}{(6-\sqrt3)(6+\sqrt3)} \\
				&= \dfrac{12\sqrt3 + 6}{36-3} \\
				&= \dfrac{12\sqrt3 + 6}{33} \\
				&= \frac4{11}\sqrt3 + \frac2{11}
			\end{align*}
		\item 
			\begin{align*}
				\dfrac{\sqrt7+6}{1+5\sqrt7} &= \dfrac{(\sqrt7+6)(5\sqrt7-1)}{(1+5\sqrt7)(5\sqrt7-1)} \\
				&= \dfrac{35-\sqrt7 + 30\sqrt7 - 6}{25\times7-1} \\
				&= \dfrac{29 + 29\sqrt7}{174} \\
				&= \frac16 + \frac16\sqrt7
			\end{align*}
	\end{enumerate}
	\end{multicols}
}

\exe{}{
	Résoudre les équations suivantes pour $x\in\R$, $x\neq0$. Exprimer $x$ sous la forme $q \sqrt{b}$ avec $q\in\Q$ rationnel et $b\in\N$ le plus petit possible à l'aide de la proposition \ref{prop:conjugué}.
	\begin{multicols}{2}
	\begin{enumerate}
		\item $\dfrac3x = \dfrac12.$
		\item $ \dfrac7{x+1}  =\dfrac12$
		\item $\dfrac8x  =\dfrac{\sqrt{3}}2$
		\item $\dfrac2{x+1} =\sqrt{3}$
		\item $\dfrac9x  =\dfrac{\sqrt{6}}5$
		\item $-\dfrac3x =3+\sqrt{7}$
		\item $-\dfrac7{2x}  =\dfrac{\sqrt{11}}3$
		\item $\dfrac3{5x+2} =3-\sqrt{5}$
	\end{enumerate}
	\end{multicols}
}{exe:équations-sqrt}{
	TODO
}


\exe{}{
	Calculer la longueur \textbf{exacte} de tous les côtés des triangles rectangles suivants à l'aide des valeurs exactes de la proposition \ref{prop:données}.

	\begin{multicols}{2}
	\includegraphics[page=4]{figures/fig-geom.pdf}
	
	\includegraphics[page=5]{figures/fig-geom.pdf}
	
	\includegraphics[page=6]{figures/fig-geom.pdf}
	
	\includegraphics[page=7]{figures/fig-geom.pdf}
	\end{multicols}
}{exe:trigo}{
	\underline{Triangle $ABC$}
	
	\begin{multicols}{2}
	\begin{align*}
		\sin(30^\circ) &= \frac{AB}{AC} \\
		\frac12 &= \frac{13}{AC} \\
		AC \frac12 &= 13 \\
		AC &= 13\times2=26
	\end{align*}
	
	\begin{align*}
		\tan(30^\circ) &= \frac{AB}{CB} \\
		\frac1{\sqrt3} &= \frac{13}{AC} \\
		AC \frac1{\sqrt3} &= 13 \\
		AC &= 13\sqrt3
	\end{align*}
	\end{multicols}
	
	\underline{Triangle $DEF$}
	
	\begin{multicols}{2}
	\begin{align*}
		\cos(60^\circ) &= \frac{EF}{DF} \\
		\frac12 &= \frac{2\sqrt6}{DF} \\
		DF \frac12 &= 2\sqrt6 \\
		DF &= 4\sqrt6
	\end{align*}
	
	\begin{align*}
		\tan(60^\circ) &= \frac{DE}{EF} \\
		\sqrt3 &= \frac{DE}{2\sqrt6} \\
		DE &= \sqrt3\times2\sqrt6 \\
		DE &= 2\sqrt{3^2\times2} = 6\sqrt2
	\end{align*}
	\end{multicols}
}


\exemulticols{}{
	Dans la figure ci-contre, calculer $x$ à l'aide du théorème de Pythagore, puis $y$ à l'aide du théorème de Thalès.
}{
	\begin{center}
		\includegraphics[page=14]{figures/fig-geom.pdf}
	\end{center}
}{exe:pyth-thalès}{
	Le théorème de Pythagore donne
		\[ (1+x)^2 = 3^2 + 3^2 = 2\times3^2. \]
	Il suit que $|1+x| = \sqrt{2\times3^2} = 3\sqrt2$, et donc
		\begin{align*}
			1+x &= 3\sqrt2 && \ou && & 1+x &= -3\sqrt2 \\
			x &= 3\sqrt2-1 && \ou && & x &= -3\sqrt2-1
		\end{align*}
	Seule la solution positive $x=3\sqrt2-1$ est à conserver.
	
	D'après le théorème de Thalès, ou en considérant le sinus de l'angle commun aux deux triangles, on obtient
		\begin{align*}
			\dfrac3{1+x} &= \dfrac{y}{x} \\
			\dfrac3{3\sqrt2} &= \dfrac{y}{3\sqrt2-1} \\
			y &= \dfrac1{\sqrt2}\cdot(3\sqrt2-1) \\
			y &= \dfrac{3\sqrt2-1}{\sqrt2} \\
			y &= \dfrac{6-\sqrt2}{2} = 3 - \frac12 \sqrt2
		\end{align*}
}


\exe{, difficulty=2}{
	En étant donné que 
		\begin{align*}
			\sin(15\degree) = \dfrac{\sqrt{6}-\sqrt{2}}4 && \text{ et } && \cos(36\degree) = \dfrac{1+\sqrt{5}}4,
		\end{align*}
	vérifier que $\cos(15\degree) = \frac{\sqrt{6}+\sqrt{2}}4$ et donner la valeur \textbf{exacte} de $\sin(36\degree)$.
}{exe:c2-s2-1}{
	Premièrement, vérifions l'égalité $\sin^2(15^\circ) + \cos^2(15^\circ) = 1$ avec les valeurs fournies.
	\begin{align*}
		\left(\frac{\sqrt{6}+\sqrt{2}}4\right)^2 + \left(\dfrac{\sqrt{6}-\sqrt{2}}4\right)^2 
		&= \frac{6+2+2\sqrt{12}}{16} + \frac{6+2-2\sqrt{12}}{16} \\
		&= \frac{8+8}{16} = 1
	\end{align*}
	Or, pour un $y$ donnée, il y a deux solutions $x$ à l'équation $x^2 + y^2 = 1$ : une positive, et l'autre négative.
	Comme $\frac{\sqrt{6}+\sqrt{2}}4 > 0$, on a trouvé la solution positive, bien égale à $\cos(15^\circ)$ car celui-ci est positif d'après le cours.
	
	Deuxièmement, $\sin(36^\circ)$ est la solution $x\in\R$ positive de
		\begin{align*}
			x^2 + \cos^2(36^\circ) &= 1 \\
			x^2 &= 1 - \left(\dfrac{1+\sqrt{5}}4\right)^2 \\
			x^2 &= 1 - \dfrac{1+5+2\sqrt5}{16} \\
			x^2 &= \dfrac{10 - 2 \sqrt5}{16} = \dfrac{5-\sqrt5}8
		\end{align*}
	Par conséquent, $\sin(36^\circ) = \sqrt{\frac{5-\sqrt5}8}$.
}

\exemulticols{}{
	Dans la figure ci-contre, calculer $z$ en sachant que $\cos(36^\circ) = \frac{1+\sqrt5}4$.
}{
	\begin{center}
		\includegraphics[page=15]{figures/fig-geom.pdf}
	\end{center}
}{exe:cos36}{
	\begin{align*}
		\cos(36^\circ) &= \dfrac{3}z \\
		\dfrac{1+\sqrt5}4 &= \dfrac3z \\
		z \cdot \dfrac{1+\sqrt5}4 &= 3 \\
		z &= 3 \cdot \dfrac4{1+\sqrt5} \\
		z &= \dfrac{12}{1+\sqrt5}  \\
		z &= \dfrac{12(\sqrt5-1)}{(1+\sqrt5)(\sqrt5 - 1)} \\
		z &=\dfrac{12(\sqrt5-1)}{4} = 3\sqrt5 - 3
	\end{align*}
}


\exe{, difficulty=1}{
	Montrer que $\sqrt2 = 1 + \frac1{1+\sqrt2}$.
	En déduire que $\sqrt2 = 1 + \frac1{2 + \frac1{1+\sqrt2}}$ et que
		\[\sqrt2 = 1 + \cfrac1{2 + \cfrac1{2 + \cfrac1{2 + \cfrac1{1+\sqrt2}}}}. \]
	Exprimer le rationnel $1 + \frac1{2 + \frac1{2 + \frac12}}$ sous la form $\frac{p}{q}$ avec $p$ et $q$ premiers entre eux et calculer $\left| \sqrt2 - \frac{p}q \right|$ au millième près.
	Faire idem avec $1 + \frac1{2 + \frac1{2 + \frac1{2 + \frac12}}}$.
}{exe:frac-continue-sqrt2}{
	En manipulant l'équation on se retrouve à vouloir montrer que $\sqrt2 - 1 = \dfrac1{\sqrt2+1}$, ce qui est équivalent à $(\sqrt2 - 1)(\sqrt2 + 1) = 1$, vrai par identité remarquable (ou par distributivité simple).
	
	Comme $\sqrt2 = 1 + \dfrac1{1+\sqrt2}$, on peut remplacer le $\sqrt2$ dans l'expression de droite par toute l'expression de droite.
	On obtient donc $\sqrt2 = 1 + \cfrac1{1 + 1 + \cfrac1{1+\sqrt2}} = 1 + \cfrac1{2 + \cfrac1{1+\sqrt2}}$ comme voulu.
	On obtient l'expression d'après en répétant le même raisonnement.
	
	Par simplifications simples, on trouve $\dfrac{p}{q} = \dfrac{17}{12}$.
	La calculatrice nous donne alors $\left| \sqrt2 - \dfrac{17}{12} \right| \approx 0,002$.
	
	Finalement, remarquons que 
	\begin{align*}
		1 + \cfrac1{2 + \cfrac1{2 + \cfrac1{2 + \cfrac12}}} &= 1 + \cfrac1{1+ 1 + \cfrac1{2 + \cfrac1{2 + \cfrac12}}}, \\
			&= 1 + \cfrac1{1+ \dfrac{17}{12}} = \dfrac{41}{29}.
	\end{align*}
	Le rationnel $\dfrac{41}{29}$ vérifie $\left| \sqrt2 - \dfrac{41}{29} \right| \approx 0,0004$, soit 5 fois plus proche de $\sqrt2$ que $\dfrac{17}{12}$ l'est.
}

\exe{, difficulty=2}{
	Le but de l'exercice est de montrer que $\sqrt2$ est \emphindex{mal approximable} par les rationnels.	
	À cette fin, pour $p, q$ deux entiers naturels non nuls, on étudie la quantité
		\[ \left|\sqrt2 - \dfrac{p}{q}\right|, \]
	distance entre $\sqrt2$ et le rationnel $\frac{p}{q}$, lorsque celle-ci est inférieure à 1.
	\begin{enumerate}
		\item Montrer que $\left|p + q \sqrt2 \right| \leq 4q$.
		\item Montrer que
			$ \left|p - q \sqrt2 \right| \cdot \left|p + q \sqrt2 \right| = \left| p^2 - 2q^2 \right| \geq 1. $
		\item En déduire que
			$  \left|\sqrt2 -  \frac{p}{q} \right| \geq \frac{1}{4q^2}. $
	\end{enumerate}
}{exe:qapprox-sqrt2}{
	TODO
}


\newpage

\section*{Exercices}
\addcontentsline{toc}{section}{Exercices}

\shipoutExercise


\section*{Approfondissements}
\addcontentsline{toc}{section}{Approfondissements}

\thm{de Dirichlet\footnotemark}{
	Pour tout nombre $x \in \R$, il existe une infinité de $p, q\in\Z$ vérifiant
		\[ \left|x - \dfrac{p}q \right| < \dfrac1{q^2}. \]
}{thm:Dirichlet}

% TODO: pourquoi la note de bas de page est pas sur la même page ? zut alors
\footnotetext{Johann Peter Gustav Lejeune Dirichlet (1805-1859), mathématicien allemand.}
%\enlargethispage{\baselineskip} 

\exe{}{
	Montrer le théorème \ref{thm:Dirichlet} dans le cas où $x\in\Q$ est un nombre rationnel.
}{exe:DirichletQ}{
	Si $x=\frac{a}b$, alors $p= k\cdot a$ et $q = k \cdot b$ pour $k=1, 2, 3, \dots$ sont une infinités de tels $p, q\in\Z$.
}

\exe{, difficulty=2}{
	Soit $x\in\R-\Q$ et $N\in\N$. Considérons $\{ 0, x, 2x, 3x, ..., Nx \}$ les $N+1$ premiers multiples de $x$.
	Montrer qu'il existe nécessairement deux entiers naturels $0\leq m<n\leq N$ tels que $|nx - mx|$ ait une partie décimale comprise entre $0$ et $\frac1N$.
	En déduire le théorème \ref{thm:Dirichlet}.
}{exe:Dirichlet}{
	On considère la partie décimale de chacun des $N+1$ multiples.
	et on découpe l'intervalle $[0;1] = [0 ; 1/N] \bigcup [1/N ; 2/N] \bigcup \cdots \bigcup [(N-1)/N ; 1]$ en $N$ sous parties de tailles $1/N$.
	
	Comme il y a $N+1$ parties décimales, mais seulement $N$ sous-intervalles, deux parties décimales appartiennent nécessairement au même sous-intervalle.
	On obtient bien $0\leq m<n\leq N$ tels que $|nx - mx|$ ait une partie décimale comprise entre $0$ et $\dfrac1N$.
	
	En factorisant par $x$, on a donc $q = n-m \leq N$ vérifiant $|qx - p| \leq \dfrac1N \leq \dfrac1q$, où $p\in\Z$ est la partie entière de $qx$.
	Voilà notre premier rationnel $p/q$ approximant $x$ avec erreur inférieure à $1/q^2$.
	
	Remarquons désormais que $|qx-p| \neq 0$ car $x \not\in\Q$, et qu'on peut choisir un $N_2$ suffisamment grand vérifiant $\dfrac1{N_2} < |qx -p|$.
	Répétons le raisonnement en remplaçant $N$ par $N_2$.
	Les entiers $p_2$ et $q_2$ trouvés sont alors nécessairement distincts des premiers car ils approximent strictement mieux $x$ : 
		\[ |q_2 x - p_2 | \leq \dfrac1{N_2} < |qx - p|. \]
	
	On prendra ensuite $N_3$ vérifiant à son tour $\dfrac1{N_3} < |q_2x -p_2|$, etc...
	En répétant le processus, on obtient bien une infinité de $p, q\in\Z$ vérifiant $|qx - p| < \dfrac1q$.
}

\nt{
	Adolf Hurwitz\footnotemark donna une version plus forte du théorème \ref{thm:Dirichlet}, avec un facteur $\dfrac1{\sqrt5}$ ajouté à droite.
	Sa démonstration utilise la théorie des \emphindex{fractions continues} dont l'exercice \ref{exe:frac-continue-sqrt2} donne un aperçu.
}

\footnotetext{Adolf Hurwitz (1859-1919), mathématicien allemand.}
%\enlargethispage{\baselineskip} 

%\newpage
\exe{, difficulty=2}{
	On pose $\Qsqrt{2} = \bigset{ r + s \sqrt2 \tq r, s \in \Q} \subseteq \R$ et on considère $x = r + s\sqrt2 \in \Qsqrt{2}$ avec $r, s \in\Q$.
	\begin{enumerate}
		\item
		Montrer que $\Q \subseteq \Qsqrt2$ est une inclusion stricte.
		\item
		Montrer que $x=0 \iff r = s = 0$.
		\item
		Montrer que si $x\neq0$, alors $\frac1x \in \Qsqrt{2}$.
		\item
		Montrer que si $y \in \Qsqrt{2}$, alors $x\cdot y \in \Qsqrt{2}$.
		\item
		Montrer que $\Qsqrt{2} \neq \R$ en montrant que $\sqrt3 \not\in \Qsqrt{2}$.
	\end{enumerate}
} {exe:Qsqrt2}{
	\begin{enumerate}
		\item
		L'inclusion $\Q \subseteq \Qsqrt2$ vient du choix de $s=0$ dans l'ensemble $\Qsqrt2$ qui permet de réaliser tous les rationnels.
		En outre, $\sqrt2 \in \Qsqrt2$ ($r=0, s=1$), et l'inclusion est donc stricte car $\sqrt2\not\in\Q$ d'après le corollaire \ref{cor:sqrt2-irrational}.
		\item
		Si $r=s=0$, alors on a immédiatement $x=0$.
		
		Réciproquement, si $x=0$, alors $r+s\sqrt2 = 0$.
		Si $s$ est non nul, on trouve $\sqrt2 = -\dfrac{r}s$, ratio de deux rationnels et donc rationnel.
		Or $\sqrt2$ est un nombre irrationnel, donc c'est impossible.
		Ainsi $s=0$ et $r=0$.
		
		\item
		Si $x\neq0$, alors 
			\[ \dfrac1x = \dfrac1{r+s\sqrt2} = \dfrac{r-s\sqrt2}{r^2 - 2s^2} = \dfrac{r}{r^2 - 2s^2} - \dfrac{s}{r^2 - 2s^2} \sqrt2 \in \Qsqrt{2}, \]
		où on a utilisé que le produit et la somme de rationnels est rationnel.
		
		\item
		Posons $y= a + b\sqrt2$ avec $a, b\in\Q$.
		Alors $x\cdot y = (r+s\sqrt2)(a+b\sqrt2) = (ra + 2sb) + (rb + sa)\sqrt2 \in \Qsqrt{2}$.
		
		\item
		Supposons par l'absurde que $\sqrt3\in\Qsqrt2$.
		Alors pour certains $r, s\in\Q$, on a $\sqrt3 = r+s\sqrt2$.
		Remarquons qu'on a $s\neq0$, car $\sqrt3$ n'est pas rationnel.
		
		En mettant au carré, on trouve $3 = r^2 + 2s^2 + 2rs\sqrt2$.
		Si $r \neq 0$, on obtient $\sqrt2 = \dfrac{3-r^2 -2s^2}{2rs} \in \Q$, ce qui n'est pas possible.
		Il suit donc que $r=0$, et donc que $\sqrt3 = s\sqrt2 \iff \sqrt6 = 2s$.
		Ainsi $\sqrt6 \in\Q$, ce qui n'est pas vrai non plus. \Large\Lightning
	\end{enumerate}
}

\nt{
	L'exercice \ref{exe:Qsqrt2} fait l'étude du corps d'extension des nombres rationnels auxquels on a ajouté la racine carrée de 2.
	La \emphindex{théorie de Galois}\footnotemark généralise cet exemple et permet notamment de répondre par la négative aux trois grands problèmes de l'Antiquité : la duplication du cube, la trisection de l'angle, et la quadrature du cercle.
}

\footnotetext{Évariste Galois (1811-1832), mathématicien français mort à l'âge de 20 ans. }

\shipoutExercise

\section*{Solutions}
\addcontentsline{toc}{section}{Solutions}

\shipoutAnswer


